\chapter{\nt{GLUT} و\nt{GABA} معًا}
\label{sec:gaba}
\begin{chaptergoals}
\item كيف يجعل الحظر $a\wedge\bar b$ من \nt{GLUT} و\nt{GABA} منظومة تامة بسلك واحد لكل بت؛
\item لماذا يقسم التثبيط التحويلي ولا يطرح، وكيف يضيّق التثبيط نافذة التزامن؛
\item كيف يختار التثبيط المتبادل مع $\beta>1$ فائزًا واحدًا، وكيف تصفّر إشارةٌ الذاكرة؛
\item متى يُبقي التثبيط السريع الشبكة مستقرة، ومتى يحوّلها التثبيط البطيء إلى إيقاع.
\end{chaptergoals}

\section{قواعد اللعبة}

شبكة عصبونات \nt{GLUT} وحدها لا تستطيع أن تختار، ولا أن تصفّر الذاكرة، ولا أن تمنع الفعالية من التصاعد المتسلسل، وكل ذلك بسبب الرتابة (القضية~\ref{prop:monotone}). لنضف إليها النوع الثاني من حيث العدد: \nt{GABA}. القواعد كما كانت: لا شيء إلا ما يحدث في الكائن الحي، وكل شيء يجري في الزمن المتصل.

نموذج العصبون هو نفسه~\eqref{eq:glutneuron}، لكن النبضة من عصبون \nt{GABA} لا ترفع جهد المستقبِل بل تخفضه. للأوزان الآن علامتان، والمرسِل يحدد العلامة، بحسب القاعدة~\eqref{eq:dalesign}.

بعض وسائط الدارة. عصبونات \nt{GABA} في القشرة من الخُمس إلى الربع~\cite{Hendry1987}. مخارجها قصيرة: تثبط جيرانها لا المناطق البعيدة. يصل التثبيط بعد جزء من الألف من الثانية ويدوم بضعة أجزاء من الألف. ومن دون مدخل يصمت عصبون \nt{GABA}: فلكي يثبط أحدًا يجب أن يتلقى هو نفسه إثارة، عادة من عصبونات \nt{GLUT} في الشبكة نفسها.

لذا يجب أن تُصنع الوصلة السالبة من عصبون \nt{GLUT} $A$ إلى العصبون $B$ عبر وسيط: $A$ يثير عصبون \nt{GABA}، وهذا يثبط $B$. ثمن قلب العلامة عصبون واحد وتأخير واحد، نحو جزء من الألف من الثانية.

في التثبيط لا يهم فقط \emph{من أين} تأتي الوصلة، بل \emph{أين} تستقر: فموضع الاستقرار يحدد إلى حد بعيد ما تفعله الوصلة~\cite{Markram2004}. مخارج \nt{GLUT} تستقر في الغالب على فروع الدخل لدى المستقبِل، أي \emph{تغصناته}، وتحمل البيانات: كل مدخل كهذا حدّ واحد في المجموع. والمخرج على جسم العصبون يؤثر في المجموع كله: هكذا يقسم كثير من عصبونات \nt{GABA} السريعة استجابة المستقبِل ويحدد متى ينطلق. والمخرج على الموضع الذي تولد فيه نبضة المستقبِل هو الكلمة الأخيرة، حظر على المخرج كله دفعة واحدة. أما المخرج على نهاية سلك عصبون آخر فيغيّر احتمال القذف $p$ لخط دخل واحد، دون المساس بغيره~\cite{Rudomin1999}؛ وهكذا تعمل، مثلًا، بوابة الألم في الفصل~\ref{sec:pain}. وخارج الدماغ تستقر المخارج على الألياف العضلية (الفصل~\ref{sec:muscles}) وعلى الأعضاء (الفصل~\ref{sec:chemoutput})، وبعضها لا يستقر على شيء بل يقذف الناقل العصبي في حجم النسيج أو في الدم (الفصلان~\ref{sec:serocomputer} و\ref{sec:chemoutput}).

\section{\foreignlanguage{english}{NOT} والحظر}

هل نستطيع الآن صنع \foreignlanguage{english}{NOT}؟ تقريبًا. العصبون الذي يعمل حين يصمت المدخل يجب أن يأخذ إثارته من مكان ما. وهي موجودة في الدماغ الحي: فعلى كل عصبون ترد باستمرار فعالية خلفية من آلاف غيره. لتبقِ الخلفية العصبون $Y$ عاملًا، وليثبطه المدخل $x$ عبر وسيط \nt{GABA}. هذه \foreignlanguage{english}{NOT}.

لكن الأنفع هو \emph{الحظر}: ”$a$، إلا عند $b$“:
\begin{empheq}[box=\keybox]{equation}
y \;=\; a \wedge \bar b .
\label{eq:veto}
\end{empheq}
العصبون $Y$ يتلقى الإثارة من $a$ والتثبيط من عصبون \nt{GABA} يثيره $b$. لا حاجة هنا إلى الخلفية: فالإثارة يعطيها $a$ نفسه. ومن الحظر و\foreignlanguage{english}{OR} يُبنى كل شيء. مثلًا \foreignlanguage{english}{XOR}: $x\oplus y=(x\wedge\bar y)\vee(\bar x\wedge y)$، أي عصبونا حظر، ووسيطا \nt{GABA}، وعصبون \foreignlanguage{english}{OR} واحد.

\begin{keyblock}
\begin{proposition}[تمام \nt{GLUT}+\nt{GABA}]
\label{prop:gabacomplete}
شبكة من عصبونات \nt{GLUT} و\nt{GABA} تستطيع أن تحسب أي دالة منطقية في المداخل، وكل بت ينتقل على سلك واحد. ولم تعد مستشعرات ”التشغيل/الإطفاء“ من الفصل~\ref{sec:glutcomputer} لازمة لذلك. وإذا وجب أن تعطي الدالة واحدًا حين تصمت كل المداخل، لزم مصدر فعالية خلفية.
\end{proposition}
\end{keyblock}

البرهان: الحظر~\eqref{eq:veto} مع \foreignlanguage{english}{OR} تامّ وظيفيًا إذا توفر واحد ثابت، لأن \foreignlanguage{english}{NOT} $x$ هو الحظر ”واحد، إلا عند $x$“. أما الدوال التي تعطي صفرًا حين تصمت كل المداخل فتُبنى من الحظر و\foreignlanguage{english}{OR} حتى من دون الواحد.

\section{اتجاه الحركة}

الحظر في الزمن، مثل \foreignlanguage{english}{AND} في الزمن في الفصل~\ref{sec:glutcomputer}، يعطي كاشفًا لاتجاه الحركة.

ليكن مستشعران متجاوران، $A$ و$B$، على بعد $\Delta x$ أحدهما من الآخر. العصبون المخرج $Y$ يتلقى الإثارة من $B$ والتثبيط من $A$ عبر وسيط \nt{GABA}، بتأخير $d$ (الشكل~\ref{fig:direction}). بقعة الضوء الجارية بسرعة $v$ من $A$ إلى $B$ تمر بـ$A$ في اللحظة $0$، فيصل التثبيط إلى $Y$ في اللحظة $d$؛ وتمر بـ$B$ في اللحظة $\Delta x/v$، وحينها تصل الإثارة. فإذا تطابقت هاتان اللحظتان ضمن النافذة~\eqref{eq:window}، أطفأ التثبيط الإثارة وصمت $Y$. ويحدث هذا عند سرعة قريبة من
\begin{empheq}[box=\keybox]{equation}
v^* \;=\; \frac{\Delta x}{d} .
\label{eq:vstar}
\end{empheq}
أما عند الحركة من $B$ إلى $A$ فتصل الإثارة من $B$ في اللحظة $0$، ولا يصل التثبيط من $A$ إلا في اللحظة $\Delta x/v+d$، وقد انطلق $Y$ عندئذ.

\begin{figure}[t]
\centering
\begin{tikzpicture}[>=Stealth,
  nrn/.style={circle,draw,very thick,fill=blue!6,minimum size=0.9cm,inner sep=0pt},
  gab/.style={nrn,fill=red!10},
  inh/.style={-{Bar[width=7pt]},thick,red!70!black}]
\node[nrn] (A) at (0,2) {$A$};
\node[nrn] (B) at (3,2) {$B$};
\node[gab] (G) at (0.9,0.6) {\small \nt{GABA}};
\node[nrn] (Y) at (3,-0.6) {$Y$};
\draw[->,thick] (A) -- (G);
\draw[inh] (G) -- node[below left=-2pt] {\scriptsize التأخير $d$} (Y);
\draw[->,thick] (B) -- (Y);
\draw[->,very thick,red!70!black] (Y) -- (4.6,-0.6) node[right,black] {\small المخرج};
\draw[->,thick,orange!85!black] (-0.6,2.9) -- (3.6,2.9) node[right,black] {\small يصمت};
\draw[<-,thick,orange!85!black] (-0.6,3.4) -- (3.6,3.4) node[right,black] {\small يستجيب};
\foreach \yy/\dd in {2.9/-1,3.4/1}{
  \foreach \k in {1,2,3}{\draw[orange!70!yellow,line width=0.8pt] (1.5+\dd*0.12+\dd*0.13*\k,\yy+0.1-0.1*\k+0.1) -- ++(\dd*0.25,0);}
  \fill[yellow!70!orange,draw=orange!70!black] (1.5,\yy) circle (0.14);}
\node[font=\scriptsize,align=center] at (1.5,3.95) {بقعة ضوء};
\end{tikzpicture}
\caption{كاشف اتجاه الحركة. $B$ يثير $Y$، و$A$ يثبطه عبر وسيط \nt{GABA} بتأخير $d$. عند الحركة من $A$ إلى $B$ بسرعة قريبة من $\Delta x/d$ يطفئ التثبيط الإثارة؛ وعند الحركة المعاكسة يتأخر. البقعة الصفراء ذات الشُّرَط هي الضوء الجاري فوق المستشعرات.}
\label{fig:direction}
\end{figure}

في الشبكية يبدو أن الانتقائية لاتجاه الحركة تنشأ بهذه الطريقة بالذات: بتثبيط غير متناظر من الجهة التي لا ينبغي أن تستدعي الحركة منها استجابة~\cite{Barlow1965}.

\section{طرح أم قسمة}

حتى الآن كان التثبيط عندنا يطرح. والحقيقة أنه أدق من ذلك.

المشبك يفتح قناة، أي يشغّل موصلية. في المشبك المثير يقع جهد الاتزان فوق العتبة بكثير، نحو $70\mV$ فوق الراحة: القناة المفتوحة تشد الجهد إلى أعلى. وفي مشبك \nt{GABA} المثبط تمرر القناة الكلور، الذي يقع جهد اتزانه قرب الراحة، فتشد القناة المفتوحة الجهد لا إلى أسفل بل \emph{نحو الراحة}. إذا قسنا الجهد من الراحة، فالجهد المستقر لغشاء بموصلية تسريب $g_L$ وموصلية مثيرة $g_E$ ومثبطة $g_I$ يساوي
\begin{empheq}[box=\keybox]{equation}
V_\infty \;=\; \frac{g_E\,E_E}{g_L+g_E+g_I} ,
\label{eq:shunt}
\end{empheq}
حيث $E_E\approx70\mV$ جهد اتزان المشبك المثير. هذا قانون أوم لمقسّم الجهد: $g_E$ تشد نحو $E_E$، و$g_L$ و$g_I$ نحو الصفر.

لا تقف $g_I$ في البسط بإشارة سالبة، بل في المقام: التثبيط لا \emph{يطرح} من الإثارة، بل \emph{يقسمها} (الشكل~\ref{fig:shunt}). يسمى هذا التثبيط تثبيطًا تحويليًا: قنوات الكلور المفتوحة تقصّر دارة الغشاء. وبالنسبة للشبكة هذا ضبط للكسب. فإذا تناسبت $g_I$ مع مجموع فعالية الجيران، استجاب كل عصبون لا للقوة المطلقة لمدخله بل لحصتها من الفعالية الكلية. وهذه القسمة على الفعالية الكلية توجد في مناطق كثيرة من الدماغ، وتُعد من عملياته المعيارية~\cite{Carandini2012}: فالشبكة نفسها تعمل مع المدخل الضعيف ومع القوي.

تحفّظ: الصيغة~\eqref{eq:shunt} تخص عصبونًا لا يطلق نبضات. أما في العصبون المنطلق فيبقى الجهد طوال الوقت قرب العتبة، ويكاد التيار عبر الموصلية المثبطة يكون ثابتًا، فيؤثر التحويل في \emph{تردد} النبضات بالطرح أساسًا~\cite{HoltKoch1997}. ويُقسَم التردد إذا أضيفت إلى العصبون في آن واحد موصليتان خلفيتان مثيرة ومثبطة متوازنتان، كما في النمط المرتعش المتوازن للقشرة الحية~\cite{Chance2002}. فالقسمة لا يحصل عليها الدماغ من التحويل وحده، بل من التثبيط مقرونًا بالضجيج الخلفي~\cite{Carandini2012}.

\begin{figure}[t]
\centering
\begin{tikzpicture}[>=Stealth,x=1.25cm,y=0.05cm]
\draw[->,gray!70!black] (0,0) -- (5.4,0) node[right,black] {\small $g_E/g_L$};
\draw[->,gray!70!black] (0,0) -- (0,68) node[above,black] {\small $V_\infty$، \foreignlanguage{english}{mV}};
\foreach \v in {20,40,60}{\draw[gray!60!black] (-0.05,\v) -- (0.05,\v) node[left=3pt,font=\scriptsize] {\v};}
\foreach \x in {1,2,3,4,5}{\draw[gray!60!black] (\x,-1.5) -- (\x,1.5) node[below=5pt,font=\scriptsize] {\x};}
\draw[very thick,blue!60!black] plot[domain=0:5,samples=60] (\x,{70*\x/(1+\x)});
\draw[very thick,red!70!black] plot[domain=0:5,samples=60] (\x,{70*\x/(3+\x)});
\draw[thick,densely dashed,gray!50!black] plot[domain=0.56:5,samples=60] (\x,{70*\x/(1+\x)-25});
\node[right,font=\small,blue!60!black] at (5.02,58.3) {بلا تثبيط};
\node[right,font=\small,red!70!black] at (5.02,45.5) {يقسم: $g_I=2g_L$};
\node[right,font=\small,gray!40!black] at (5.02,32.5) {يطرح $25\mV$};
\end{tikzpicture}
\caption{التثبيط التحويلي يقسم الجهد ولا يطرح منه. المنحنى الأزرق هو الجهد المستقر~\eqref{eq:shunt} بلا تثبيط، والأحمر مع موصلية مثبطة $g_I=2g_L$. الأحمر هو الأزرق نفسه ممطوطًا أفقيًا ثلاث مرات: كأن المدخل قُسم على $1+g_I/g_L=3$. والمتقطع تثبيط يطرح $25\mV$ ثابتة: ينخفض المنحنى كله ولا يتغير الميل.}
\label{fig:shunt}
\end{figure}

وثمة نتيجة ثانية. ثابت زمن الغشاء $\tau_{\text{eff}}=C/(g_L+g_E+g_I)$، والتثبيط يقصّره. ونافذة التزامن~\eqref{eq:window} متناسبة مع $\tau$، فالتثبيط الذي يصل مع الإثارة \emph{يضيّق النافذة}. والدارة التي يثير فيها المدخلُ العصبونَ مباشرة ويثبطه في الوقت نفسه، بتأخير جزء من الألف من الثانية، عبر وسيط \nt{GABA}، من أكثر الدارات شيوعًا في الدماغ: لا يلحق العصبون أن يستجيب إلا لما وصل في الجزء الأول أو الثاني من الألف من الثانية~\cite{Pouille2001}.

\section{الاختيار}

والآن الأهم مما كان ينقص شبكة \nt{GLUT}: اختيار واحد من عدة. لنأخذ مجموعتين من عصبونات \nt{GLUT} بمدخلين $I_1$ و$I_2$. كل منهما تثبط الأخرى عبر وسطاء \nt{GABA} الخاصين بها بقوة $\beta$ (الشكل~\ref{fig:wta}). وللترددين $r_1,r_2$ نحصل على
\begin{equation}
\tau\,\frac{\dd r_1}{\dd t} = -r_1 + \bigl[I_1-\beta r_2\bigr]_+ ,\qquad
\tau\,\frac{\dd r_2}{\dd t} = -r_2 + \bigl[I_2-\beta r_1\bigr]_+ ,
\label{eq:wta}
\end{equation}
حيث $[u]_+=\max(u,0)$: لا يكون التردد سالبًا.

\begin{figure}[t]
\centering
\begin{tikzpicture}[>=Stealth,
  nrn/.style={circle,draw,very thick,fill=blue!6,minimum size=0.9cm,inner sep=0pt},
  gab/.style={nrn,fill=red!10,minimum size=0.75cm},
  inh/.style={-{Bar[width=7pt]},thick,red!70!black}]
\node[nrn] (E1) at (0,0) {$r_1$};
\node[nrn] (E2) at (4,0) {$r_2$};
\node[gab] (G1) at (1.4,1.1) {};
\node[gab] (G2) at (2.6,-1.1) {};
\draw[->,thick] (E1) -- (G1); \draw[inh] (G1) -- (E2);
\draw[->,thick] (E2) -- (G2); \draw[inh] (G2) -- (E1);
\draw[->,thick,blue!60!black] (-1.6,0) node[left,black] {$I_1$} -- (E1);
\draw[->,thick,blue!60!black] (5.6,0) node[right,black] {$I_2$} -- (E2);
\end{tikzpicture}
\caption{الاختيار من اثنين. مجموعتان من عصبونات \nt{GLUT} (بالأزرق) تثبط إحداهما الأخرى عبر وسطاء \nt{GABA} (بالأحمر).}
\label{fig:wta}
\end{figure}

ما دام العصبونان كلاهما يعملان، فالنظام~\eqref{eq:wta} خطي، ويحدد سلوكه قيمتان ذاتيتان: $-(1+\beta)/\tau$ للانحرافات المتماثلة في الترددين، و$-(1-\beta)/\tau$ للمتعاكسة. عند التثبيط الضعيف، $\beta<1$، القيمتان سالبتان: يعمل العصبونان كلاهما، والتثبيط لا يفعل سوى أن يخفت الأضعف. وعند التثبيط القوي، $\beta>1$، تصبح القيمة الثانية موجبة: أي فرق بين $r_1$ و$r_2$ ينمو حتى يصمت أحد العصبونين تمامًا.

\begin{keyblock}
\begin{proposition}[الفائز يأخذ كل شيء]
\label{prop:wta}
إذا كان التثبيط المتبادل أقوى من الواحد، $\beta>1$، فالحالة التي تعمل فيها المجموعتان كلتاهما غير مستقرة. تصل الشبكة إلى حالة تعمل فيها مجموعة واحدة وتصمت الأخرى. والفائز بالتردد $r_1=I_1$ يحتفظ بفوزه ما دام $I_2<\beta I_1$.
\end{proposition}
\end{keyblock}

اختبرنا ذلك على النموذج. عند $\beta=0.5$ والمدخلين $1.0$ و$0.9$ تعمل المجموعتان كلتاهما بترددين $0.73$ و$0.53$. وعند $\beta=2$ والمدخلين أنفسهما تصمت المجموعة الثانية تمامًا. ولهذا الاختيار ذاكرة: إذا تبادل المدخلان مكانيهما لاحقًا، بقي الفائز السابق فائزًا، لأن $1.0<2\cdot0.9$.

ومع مئة مجموعة يبقى الأمر كذلك: كل منها تثبط البقية، وتبقى واحدة.

\section{تصفير الذاكرة}

ذاكرة الفصل~\ref{sec:glutcomputer} لم تكن تنطفئ إلا بالتعب. لنضف إليها عصبون \nt{GABA} تثيره إشارة ”تصفير“ ويثبط المجموعة. يخفض التثبيط المنحنى $f(wr+I)$ في الشكل~\ref{fig:bistable}، فيختفي التقاطع العلوي، وتنطفئ المجموعة، وتبقى في الحالة السفلى بعد انتهاء التصفير. الآن تُكتب الذاكرة بإشارة وتُمحى بأخرى، كما في قلّاب بمدخلي ضبط وتصفير.

والأجمل أن نصل مجموعتين كهاتين بتثبيط متبادل، كما في الشكل~\ref{fig:wta}، ونعطي كلًّا منهما تغذية راجعة على نفسها. الإشارة على مدخل إحدى المجموعتين تنقل الخلية إلى حالتها، وتتذكر الخلية القرار الأخير. هذا نظير مباشر للماسك المؤلف من عنصري \foreignlanguage{english}{NOR} موصولين تقاطعيًا.

\section{الاستقرار والإيقاع}

تبقى العلة الثالثة لشبكة \nt{GLUT}: الانهيار المتسلسل. لنأخذ مجموعة من عصبونات \nt{GLUT} بتغذية راجعة على نفسها $w_{EE}$، ومجموعة من عصبونات \nt{GABA} تثيرها الأولى بقوة $w_{IE}$ وتثبط الأولى بقوة $w_{EI}$. وللترددين $E$ و$I$ نحصل على
\begin{equation}
\tau_E\frac{\dd E}{\dd t} = -E + w_{EE}E - w_{EI}I + h, \qquad
\tau_I\frac{\dd I}{\dd t} = -I + w_{IE}E ,
\label{eq:ei}
\end{equation}
حيث $h$ المدخل الخارجي~\cite{WilsonCowan1972}. من دون تثبيط، عند $w_{EE}>1$، تنمو الفعالية بلا حد: كل نبضة تولّد أكثر من نبضة لاحقة. ومع التثبيط يكون التردد المستقر
\begin{equation}
E^* \;=\; \frac{h}{1-w_{EE}+w_{EI}w_{IE}}
\label{eq:eistar}
\end{equation}
منتهيًا إذا كان المقام موجبًا. لكن هذا لا يكفي: يجب أيضًا أن يكون الاتزان جاذبًا. ومن التحليل الخطي لـ~\eqref{eq:ei} نحصل على شرطين.

\begin{keyblock}
\begin{proposition}[شروط الاستقرار]
\label{prop:ei}
الاتزان~\eqref{eq:eistar} مستقر إذا كان التثبيط
\begin{enumerate}
\item[(i)] \emph{قويًا بما يكفي}: $w_{EI}\,w_{IE} > w_{EE}-1$؛
\item[(ii)] \emph{سريعًا بما يكفي}: $\tau_I < \tau_E/(w_{EE}-1)$.
\end{enumerate}
إذا انتُهك (i) تصاعدت الفعالية تصاعدًا متسلسلًا. وإذا تحقق (i) وانتُهك (ii) تأخر التثبيط، وبدأت الفعالية تتذبذب.
\end{proposition}
\end{keyblock}

الشرط~(i) هو محدد مصفوفة النظام~\eqref{eq:ei}، والشرط~(ii) أثرها. ومعنى الثاني واضح لكل من ضبط منظِّمًا: التغذية الراجعة المتأخرة لا تخمد الانحراف بل تؤرجحه.

\begin{figure}[t]
\centering
\begin{tikzpicture}[>=Stealth]
\draw[->,gray!70!black] (0,0) -- (10.9,0) node[right,black] {\small $t$، \foreignlanguage{english}{ms}};
\draw[->,gray!70!black] (0,0) -- (0,3.3) node[left,black] {\small $E$};
\foreach \t in {0,50,100,150}{\draw[gray!70!black] (\t*0.07,0.06) -- (\t*0.07,-0.06) node[below,font=\scriptsize] {\t};}
\draw[very thick,gray!60!black,densely dashed] plot coordinates {(0.000,0.000) (0.035,0.071) (0.070,0.144) (0.105,0.218) (0.140,0.294) (0.175,0.373) (0.210,0.453) (0.245,0.535) (0.280,0.620) (0.315,0.706) (0.350,0.795) (0.385,0.886) (0.420,0.979) (0.455,1.075) (0.490,1.173) (0.525,1.274) (0.560,1.377) (0.595,1.482) (0.630,1.591) (0.665,1.702) (0.700,1.816) (0.735,1.933) (0.770,2.052) (0.805,2.175) (0.840,2.301) (0.875,2.430) (0.910,2.563) (0.945,2.698) (0.980,2.838) (1.015,2.980) (1.050,3.080) (1.085,3.080) (1.120,3.080) (1.155,3.080) (1.190,3.080) (1.225,3.080) (1.260,3.080) (1.295,3.080) (1.330,3.080) (1.365,3.080) (1.400,3.080) (1.435,3.080) (1.470,3.080) (1.505,3.080) (1.540,3.080) (1.575,3.080) (1.610,3.080) (1.645,3.080) (1.680,3.080) (1.715,3.080) (1.750,3.080) (1.785,3.080) (1.820,3.080) (1.855,3.080) (1.890,3.080) (1.925,3.080) (1.960,3.080) (1.995,3.080) (2.030,3.080) (2.065,3.080) (2.100,3.080) (2.135,3.080) (2.170,3.080) (2.205,3.080) (2.240,3.080) (2.275,3.080) (2.310,3.080) (2.345,3.080) (2.380,3.080) (2.415,3.080) (2.450,3.080) (2.485,3.080) (2.520,3.080) (2.555,3.080) (2.590,3.080) (2.625,3.080) (2.660,3.080) (2.695,3.080) (2.730,3.080) (2.765,3.080) (2.800,3.080) (2.835,3.080) (2.870,3.080) (2.905,3.080) (2.940,3.080) (2.975,3.080) (3.010,3.080) (3.045,3.080) (3.080,3.080) (3.115,3.080) (3.150,3.080) (3.185,3.080) (3.220,3.080) (3.255,3.080) (3.290,3.080) (3.325,3.080) (3.360,3.080) (3.395,3.080) (3.430,3.080) (3.465,3.080) (3.500,3.080) (3.535,3.080) (3.570,3.080) (3.605,3.080) (3.640,3.080) (3.675,3.080) (3.710,3.080) (3.745,3.080) (3.780,3.080) (3.815,3.080) (3.850,3.080) (3.885,3.080) (3.920,3.080) (3.955,3.080) (3.990,3.080) (4.025,3.080) (4.060,3.080) (4.095,3.080) (4.130,3.080) (4.165,3.080) (4.200,3.080) (4.235,3.080) (4.270,3.080) (4.305,3.080) (4.340,3.080) (4.375,3.080) (4.410,3.080) (4.445,3.080) (4.480,3.080) (4.515,3.080) (4.550,3.080) (4.585,3.080) (4.620,3.080) (4.655,3.080) (4.690,3.080) (4.725,3.080) (4.760,3.080) (4.795,3.080) (4.830,3.080) (4.865,3.080) (4.900,3.080) (4.935,3.080) (4.970,3.080) (5.005,3.080) (5.040,3.080) (5.075,3.080) (5.110,3.080) (5.145,3.080) (5.180,3.080) (5.215,3.080) (5.250,3.080) (5.285,3.080) (5.320,3.080) (5.355,3.080) (5.390,3.080) (5.425,3.080) (5.460,3.080) (5.495,3.080) (5.530,3.080) (5.565,3.080) (5.600,3.080) (5.635,3.080) (5.670,3.080) (5.705,3.080) (5.740,3.080) (5.775,3.080) (5.810,3.080) (5.845,3.080) (5.880,3.080) (5.915,3.080) (5.950,3.080) (5.985,3.080) (6.020,3.080) (6.055,3.080) (6.090,3.080) (6.125,3.080) (6.160,3.080) (6.195,3.080) (6.230,3.080) (6.265,3.080) (6.300,3.080) (6.335,3.080) (6.370,3.080) (6.405,3.080) (6.440,3.080) (6.475,3.080) (6.510,3.080) (6.545,3.080) (6.580,3.080) (6.615,3.080) (6.650,3.080) (6.685,3.080) (6.720,3.080) (6.755,3.080) (6.790,3.080) (6.825,3.080) (6.860,3.080) (6.895,3.080) (6.930,3.080) (6.965,3.080) (7.000,3.080) (7.035,3.080) (7.070,3.080) (7.105,3.080) (7.140,3.080) (7.175,3.080) (7.210,3.080) (7.245,3.080) (7.280,3.080) (7.315,3.080) (7.350,3.080) (7.385,3.080) (7.420,3.080) (7.455,3.080) (7.490,3.080) (7.525,3.080) (7.560,3.080) (7.595,3.080) (7.630,3.080) (7.665,3.080) (7.700,3.080) (7.735,3.080) (7.770,3.080) (7.805,3.080) (7.840,3.080) (7.875,3.080) (7.910,3.080) (7.945,3.080) (7.980,3.080) (8.015,3.080) (8.050,3.080) (8.085,3.080) (8.120,3.080) (8.155,3.080) (8.190,3.080) (8.225,3.080) (8.260,3.080) (8.295,3.080) (8.330,3.080) (8.365,3.080) (8.400,3.080) (8.435,3.080) (8.470,3.080) (8.505,3.080) (8.540,3.080) (8.575,3.080) (8.610,3.080) (8.645,3.080) (8.680,3.080) (8.715,3.080) (8.750,3.080) (8.785,3.080) (8.820,3.080) (8.855,3.080) (8.890,3.080) (8.925,3.080) (8.960,3.080) (8.995,3.080) (9.030,3.080) (9.065,3.080) (9.100,3.080) (9.135,3.080) (9.170,3.080) (9.205,3.080) (9.240,3.080) (9.275,3.080) (9.310,3.080) (9.345,3.080) (9.380,3.080) (9.415,3.080) (9.450,3.080) (9.485,3.080) (9.520,3.080) (9.555,3.080) (9.590,3.080) (9.625,3.080) (9.660,3.080) (9.695,3.080) (9.730,3.080) (9.765,3.080) (9.800,3.080) (9.835,3.080) (9.870,3.080) (9.905,3.080) (9.940,3.080) (9.975,3.080) (10.010,3.080) (10.045,3.080) (10.080,3.080) (10.115,3.080) (10.150,3.080) (10.185,3.080) (10.220,3.080) (10.255,3.080) (10.290,3.080) (10.325,3.080) (10.360,3.080) (10.395,3.080) (10.430,3.080) (10.465,3.080) (10.500,3.080)};
\draw[very thick,blue!60!black] plot coordinates {(0.000,0.000) (0.035,0.071) (0.070,0.141) (0.105,0.211) (0.140,0.279) (0.175,0.344) (0.210,0.406) (0.245,0.465) (0.280,0.519) (0.315,0.570) (0.350,0.616) (0.385,0.658) (0.420,0.697) (0.455,0.731) (0.490,0.763) (0.525,0.790) (0.560,0.815) (0.595,0.836) (0.630,0.855) (0.665,0.871) (0.700,0.886) (0.735,0.898) (0.770,0.908) (0.805,0.916) (0.840,0.924) (0.875,0.929) (0.910,0.934) (0.945,0.938) (0.980,0.941) (1.015,0.943) (1.050,0.945) (1.085,0.946) (1.120,0.946) (1.155,0.947) (1.190,0.947) (1.225,0.947) (1.260,0.946) (1.295,0.946) (1.330,0.945) (1.365,0.944) (1.400,0.944) (1.435,0.943) (1.470,0.942) (1.505,0.941) (1.540,0.941) (1.575,0.940) (1.610,0.939) (1.645,0.939) (1.680,0.938) (1.715,0.937) (1.750,0.937) (1.785,0.936) (1.820,0.936) (1.855,0.936) (1.890,0.935) (1.925,0.935) (1.960,0.935) (1.995,0.934) (2.030,0.934) (2.065,0.934) (2.100,0.934) (2.135,0.934) (2.170,0.934) (2.205,0.934) (2.240,0.933) (2.275,0.933) (2.310,0.933) (2.345,0.933) (2.380,0.933) (2.415,0.933) (2.450,0.933) (2.485,0.933) (2.520,0.933) (2.555,0.933) (2.590,0.933) (2.625,0.933) (2.660,0.933) (2.695,0.933) (2.730,0.933) (2.765,0.933) (2.800,0.933) (2.835,0.933) (2.870,0.933) (2.905,0.933) (2.940,0.933) (2.975,0.933) (3.010,0.933) (3.045,0.933) (3.080,0.933) (3.115,0.933) (3.150,0.933) (3.185,0.933) (3.220,0.933) (3.255,0.933) (3.290,0.933) (3.325,0.933) (3.360,0.933) (3.395,0.933) (3.430,0.933) (3.465,0.933) (3.500,0.933) (3.535,0.933) (3.570,0.933) (3.605,0.933) (3.640,0.933) (3.675,0.933) (3.710,0.933) (3.745,0.933) (3.780,0.933) (3.815,0.933) (3.850,0.933) (3.885,0.933) (3.920,0.933) (3.955,0.933) (3.990,0.933) (4.025,0.933) (4.060,0.933) (4.095,0.933) (4.130,0.933) (4.165,0.933) (4.200,0.933) (4.235,0.933) (4.270,0.933) (4.305,0.933) (4.340,0.933) (4.375,0.933) (4.410,0.933) (4.445,0.933) (4.480,0.933) (4.515,0.933) (4.550,0.933) (4.585,0.933) (4.620,0.933) (4.655,0.933) (4.690,0.933) (4.725,0.933) (4.760,0.933) (4.795,0.933) (4.830,0.933) (4.865,0.933) (4.900,0.933) (4.935,0.933) (4.970,0.933) (5.005,0.933) (5.040,0.933) (5.075,0.933) (5.110,0.933) (5.145,0.933) (5.180,0.933) (5.215,0.933) (5.250,0.933) (5.285,0.933) (5.320,0.933) (5.355,0.933) (5.390,0.933) (5.425,0.933) (5.460,0.933) (5.495,0.933) (5.530,0.933) (5.565,0.933) (5.600,0.933) (5.635,0.933) (5.670,0.933) (5.705,0.933) (5.740,0.933) (5.775,0.933) (5.810,0.933) (5.845,0.933) (5.880,0.933) (5.915,0.933) (5.950,0.933) (5.985,0.933) (6.020,0.933) (6.055,0.933) (6.090,0.933) (6.125,0.933) (6.160,0.933) (6.195,0.933) (6.230,0.933) (6.265,0.933) (6.300,0.933) (6.335,0.933) (6.370,0.933) (6.405,0.933) (6.440,0.933) (6.475,0.933) (6.510,0.933) (6.545,0.933) (6.580,0.933) (6.615,0.933) (6.650,0.933) (6.685,0.933) (6.720,0.933) (6.755,0.933) (6.790,0.933) (6.825,0.933) (6.860,0.933) (6.895,0.933) (6.930,0.933) (6.965,0.933) (7.000,0.933) (7.035,0.933) (7.070,0.933) (7.105,0.933) (7.140,0.933) (7.175,0.933) (7.210,0.933) (7.245,0.933) (7.280,0.933) (7.315,0.933) (7.350,0.933) (7.385,0.933) (7.420,0.933) (7.455,0.933) (7.490,0.933) (7.525,0.933) (7.560,0.933) (7.595,0.933) (7.630,0.933) (7.665,0.933) (7.700,0.933) (7.735,0.933) (7.770,0.933) (7.805,0.933) (7.840,0.933) (7.875,0.933) (7.910,0.933) (7.945,0.933) (7.980,0.933) (8.015,0.933) (8.050,0.933) (8.085,0.933) (8.120,0.933) (8.155,0.933) (8.190,0.933) (8.225,0.933) (8.260,0.933) (8.295,0.933) (8.330,0.933) (8.365,0.933) (8.400,0.933) (8.435,0.933) (8.470,0.933) (8.505,0.933) (8.540,0.933) (8.575,0.933) (8.610,0.933) (8.645,0.933) (8.680,0.933) (8.715,0.933) (8.750,0.933) (8.785,0.933) (8.820,0.933) (8.855,0.933) (8.890,0.933) (8.925,0.933) (8.960,0.933) (8.995,0.933) (9.030,0.933) (9.065,0.933) (9.100,0.933) (9.135,0.933) (9.170,0.933) (9.205,0.933) (9.240,0.933) (9.275,0.933) (9.310,0.933) (9.345,0.933) (9.380,0.933) (9.415,0.933) (9.450,0.933) (9.485,0.933) (9.520,0.933) (9.555,0.933) (9.590,0.933) (9.625,0.933) (9.660,0.933) (9.695,0.933) (9.730,0.933) (9.765,0.933) (9.800,0.933) (9.835,0.933) (9.870,0.933) (9.905,0.933) (9.940,0.933) (9.975,0.933) (10.010,0.933) (10.045,0.933) (10.080,0.933) (10.115,0.933) (10.150,0.933) (10.185,0.933) (10.220,0.933) (10.255,0.933) (10.290,0.933) (10.325,0.933) (10.360,0.933) (10.395,0.933) (10.430,0.933) (10.465,0.933) (10.500,0.933)};
\draw[very thick,red!70!black] plot coordinates {(0.000,0.000) (0.035,0.071) (0.070,0.143) (0.105,0.217) (0.140,0.293) (0.175,0.370) (0.210,0.449) (0.245,0.528) (0.280,0.609) (0.315,0.691) (0.350,0.774) (0.385,0.858) (0.420,0.943) (0.455,1.028) (0.490,1.114) (0.525,1.200) (0.560,1.287) (0.595,1.374) (0.630,1.462) (0.665,1.549) (0.700,1.636) (0.735,1.724) (0.770,1.810) (0.805,1.897) (0.840,1.983) (0.875,2.068) (0.910,2.153) (0.945,2.237) (0.980,2.320) (1.015,2.402) (1.050,2.483) (1.085,2.563) (1.120,2.641) (1.155,2.717) (1.190,2.792) (1.225,2.866) (1.260,2.937) (1.295,3.007) (1.330,3.075) (1.365,3.080) (1.400,3.080) (1.435,3.080) (1.470,3.080) (1.505,3.080) (1.540,3.080) (1.575,3.080) (1.610,3.080) (1.645,3.080) (1.680,3.080) (1.715,3.080) (1.750,3.080) (1.785,3.080) (1.820,3.080) (1.855,3.080) (1.890,3.080) (1.925,3.080) (1.960,3.080) (1.995,3.080) (2.030,3.080) (2.065,3.080) (2.100,3.080) (2.135,3.080) (2.170,3.080) (2.205,3.080) (2.240,3.080) (2.275,3.080) (2.310,3.080) (2.345,3.080) (2.380,3.080) (2.415,3.080) (2.450,3.080) (2.485,3.080) (2.520,3.080) (2.555,3.080) (2.590,3.080) (2.625,3.080) (2.660,3.080) (2.695,3.080) (2.730,3.080) (2.765,3.080) (2.800,3.080) (2.835,3.052) (2.870,2.973) (2.905,2.892) (2.940,2.807) (2.975,2.719) (3.010,2.629) (3.045,2.535) (3.080,2.438) (3.115,2.339) (3.150,2.238) (3.185,2.133) (3.220,2.029) (3.255,1.930) (3.290,1.836) (3.325,1.747) (3.360,1.661) (3.395,1.580) (3.430,1.503) (3.465,1.430) (3.500,1.360) (3.535,1.294) (3.570,1.231) (3.605,1.171) (3.640,1.113) (3.675,1.059) (3.710,1.007) (3.745,0.958) (3.780,0.911) (3.815,0.867) (3.850,0.825) (3.885,0.784) (3.920,0.746) (3.955,0.710) (3.990,0.675) (4.025,0.642) (4.060,0.611) (4.095,0.581) (4.130,0.553) (4.165,0.526) (4.200,0.500) (4.235,0.476) (4.270,0.452) (4.305,0.430) (4.340,0.409) (4.375,0.389) (4.410,0.370) (4.445,0.352) (4.480,0.335) (4.515,0.319) (4.550,0.303) (4.585,0.288) (4.620,0.274) (4.655,0.261) (4.690,0.248) (4.725,0.236) (4.760,0.225) (4.795,0.214) (4.830,0.203) (4.865,0.193) (4.900,0.184) (4.935,0.175) (4.970,0.166) (5.005,0.158) (5.040,0.151) (5.075,0.143) (5.110,0.136) (5.145,0.130) (5.180,0.123) (5.215,0.117) (5.250,0.111) (5.285,0.106) (5.320,0.101) (5.355,0.096) (5.390,0.091) (5.425,0.087) (5.460,0.083) (5.495,0.079) (5.530,0.075) (5.565,0.071) (5.600,0.068) (5.635,0.064) (5.670,0.061) (5.705,0.058) (5.740,0.055) (5.775,0.053) (5.810,0.050) (5.845,0.048) (5.880,0.045) (5.915,0.043) (5.950,0.041) (5.985,0.039) (6.020,0.037) (6.055,0.035) (6.090,0.034) (6.125,0.032) (6.160,0.030) (6.195,0.029) (6.230,0.027) (6.265,0.026) (6.300,0.025) (6.335,0.024) (6.370,0.022) (6.405,0.021) (6.440,0.021) (6.475,0.022) (6.510,0.024) (6.545,0.027) (6.580,0.032) (6.615,0.037) (6.650,0.044) (6.685,0.052) (6.720,0.061) (6.755,0.071) (6.790,0.083) (6.825,0.095) (6.860,0.109) (6.895,0.124) (6.930,0.140) (6.965,0.157) (7.000,0.176) (7.035,0.195) (7.070,0.215) (7.105,0.237) (7.140,0.260) (7.175,0.283) (7.210,0.308) (7.245,0.333) (7.280,0.360) (7.315,0.388) (7.350,0.416) (7.385,0.445) (7.420,0.476) (7.455,0.507) (7.490,0.539) (7.525,0.571) (7.560,0.604) (7.595,0.638) (7.630,0.673) (7.665,0.708) (7.700,0.744) (7.735,0.781) (7.770,0.818) (7.805,0.855) (7.840,0.893) (7.875,0.931) (7.910,0.969) (7.945,1.008) (7.980,1.047) (8.015,1.086) (8.050,1.126) (8.085,1.165) (8.120,1.204) (8.155,1.244) (8.190,1.283) (8.225,1.322) (8.260,1.362) (8.295,1.400) (8.330,1.439) (8.365,1.477) (8.400,1.515) (8.435,1.553) (8.470,1.590) (8.505,1.626) (8.540,1.662) (8.575,1.698) (8.610,1.733) (8.645,1.767) (8.680,1.800) (8.715,1.832) (8.750,1.864) (8.785,1.895) (8.820,1.924) (8.855,1.953) (8.890,1.981) (8.925,2.007) (8.960,2.033) (8.995,2.057) (9.030,2.080) (9.065,2.102) (9.100,2.123) (9.135,2.142) (9.170,2.160) (9.205,2.177) (9.240,2.192) (9.275,2.205) (9.310,2.217) (9.345,2.228) (9.380,2.237) (9.415,2.245) (9.450,2.251) (9.485,2.255) (9.520,2.258) (9.555,2.259) (9.590,2.259) (9.625,2.256) (9.660,2.253) (9.695,2.247) (9.730,2.240) (9.765,2.231) (9.800,2.220) (9.835,2.208) (9.870,2.194) (9.905,2.178) (9.940,2.160) (9.975,2.141) (10.010,2.120) (10.045,2.098) (10.080,2.074) (10.115,2.048) (10.150,2.021) (10.185,1.992) (10.220,1.961) (10.255,1.929) (10.290,1.895) (10.325,1.860) (10.360,1.824) (10.395,1.786) (10.430,1.746) (10.465,1.706) (10.500,1.663)};

\node[font=\small,gray!40!black,anchor=west] at (1.3,3.45) {بلا تثبيط: انهيار متسلسل};
\node[font=\small,blue!60!black,anchor=west] at (7.4,1.25) {تثبيط سريع};
\node[font=\small,red!70!black,anchor=west] at (7.4,2.55) {تثبيط بطيء};
\end{tikzpicture}
\caption{تردد مجموعة \nt{GLUT} ذات تغذية راجعة $w_{EE}=1.5$ و$\tau_E=10\ms$، بعد تشغيل المدخل. بلا تثبيط (المتقطع) تنمو الفعالية بلا حد. ومع تثبيط بقوة $w_{EI}w_{IE}=2$ و$\tau_I=2\ms$ (المنحنى الأزرق) تستقر عند المستوى $E^*=h/(1-1.5+2)=0.67h$. ومع التثبيط نفسه لكن بطيئًا، $\tau_I=30\ms$ (المنحنى الأحمر)، تتذبذب الفعالية.}
\label{fig:ei}
\end{figure}

يُعتقد أن القشرة تعمل في هذا النمط بالذات: تغذيتها الراجعة الذاتية قوية بما يكفي لتنجرف إلى انهيار متسلسل لولا التثبيط، ولا يبقيها عند نقطة التشغيل إلا التثبيط السريع~\cite{Tsodyks1997isn}.

ولهذا النمط بصمة متناقضة ظاهريًا يمكن التعرف عليه بها من الخارج~\cite{Tsodyks1997isn}. لنضف إلى~\eqref{eq:ei} مدخلًا خارجيًا $h_I$ مباشرة إلى مجموعة \nt{GABA}. تصبح الترددات المستقرة
\begin{equation}
E^* \;=\; \frac{h-w_{EI}\,h_I}{D},\qquad
I^* \;=\; \frac{w_{IE}\,h+(1-w_{EE})\,h_I}{D},\qquad
D \;=\; 1-w_{EE}+w_{EI}w_{IE} .
\label{eq:isn}
\end{equation}
إذا كان $w_{EE}>1$ فمعامل $h_I$ في الصيغة الثانية سالب: ادفع عصبونات \nt{GABA} فـ\emph{يهبط} ترددها. والسبب في الحلقة. التثبيط الزائد يخفت مجموعة \nt{GLUT}، فتقل إثارتها لعصبونات \nt{GABA}، فتخسر هذه أكثر مما كسبت. وبأرقام الشكل~\ref{fig:ei} ($w_{EE}=1.5$، $w_{EI}w_{IE}=2$، $D=1.5$) كل وحدة من المدخل المضاف تخفض $I^*$ بثلث وحدة. وقد وُجدت هذه البصمة في القشرة: حين يُكبت استجابة عصبونات القشرة البصرية منبهٌ محيط، فإن التيار المثبط الذي تتلقاه لا يزداد بل يهبط مع المثير~\cite{Ozeki2009}. ولاحقًا وُجدت البصمة نفسها في مناطق وطبقات مختلفة من القشرة~\cite{Sanzeni2020}.

والتذبذبات عند انتهاك الشرط~(ii) موجودة في الدماغ أيضًا. الحلقة ”\nt{GLUT} يثير \nt{GABA}، و\nt{GABA} يثبط \nt{GLUT}“ بتأخير بضعة أجزاء من الألف من الثانية تولّد إيقاعًا دوره من رتبة $12\text{--}50\ms$ (التردد $20\text{--}80\,\text{Hz}$)، وهذه الإيقاعات السريعة في القشرة معروفة جيدًا~\cite{Cardin2009,Buzsaki2012}. هذه ليست ساعة الحاسوب: الإيقاع محلي ولا ينشأ إلا حين تكون الشبكة نشطة. لكنه على مدى بضع دورات يقسم الزمن إلى نوافذ يمكن للعصبونات أن تنطلق فيها.

\section{الخلاصة}

\begin{table}[t]
\centering
\small
\begin{tabular}{>{\raggedright}p{3.3cm}>{\raggedright}p{4.4cm}>{\raggedright\arraybackslash}p{4.4cm}}
\toprule
& \nt{GLUT} فقط & \nt{GLUT} + \nt{GABA} \\
\midrule
\foreignlanguage{english}{NOT} & عبر مستشعرات ”التشغيل/الإطفاء“ فقط & الحظر $a\wedge\bar b$؛ و\foreignlanguage{english}{NOT} مع فعالية خلفية \\
الأسلاك لكل بت & اثنان & واحد \\
اتجاه الحركة & لا & حظر مع تأخير \\
ضبط الكسب & لا & القسمة: التحويل مع الضجيج الخلفي \\
نافذة التزامن & $\sim\tau$ & يضيّقها التثبيط \\
اختيار واحد من كثيرين & لا & تثبيط متبادل، $\beta>1$ \\
تصفير الذاكرة & بالتعب فقط & بإشارة عبر \nt{GABA} \\
الاستقرار & بالتعب فقط & تغذية راجعة سالبة سريعة \\
الإيقاع & غير لازم & ينشأ مع التثبيط البطيء \\
\bottomrule
\end{tabular}
\caption{ما يضيفه \nt{GABA} إلى شبكة من عصبونات \nt{GLUT}.}
\label{tab:gaba}
\end{table}

لا يكاد \nt{GABA} يضيف إلى الشبكة \emph{حسابات} جديدة، فكل ذلك كان يمكن حسابه بالمستشعرات ثنائية السلك، بل يضيف \emph{التحكم} (الجدول~\ref{tab:gaba}): الاختيار، والتصفير، وضبط الكسب، والاستقرار. الإثارة في الدماغ تحمل البيانات، والتثبيط يقرر أي البيانات تمر، ومتى، وبأي علوّ. ولعصبون من كل أربعة أو خمسة في القشرة هذا تقسيم معقول جدًا للعمل.

\section{تمارين}

\begin{exercise}[\lvlA]
\label{ex:03:1}
\emph{التحويل بالأرقام.} $E_E=70\mV$. بحسب~\eqref{eq:shunt} أوجد $V_\infty$ عند $g_E=g_L$ و$4g_L$، بلا تثبيط ومع تحويل $g_I=2g_L$. أي $V_\infty$ عند $g_E=4g_L$ يعطيه طرحٌ يُنقص عند $g_E=g_L$ قدر ما ينقصه التحويل؟ كم مرة يضيّق التحويل نافذة التزامن؟
\end{exercise}

\begin{exercise}[\lvlB]
\label{ex:03:2}
\emph{استنتاج شروط الاستقرار.} أوجد أثر مصفوفة النظام الخطي~\eqref{eq:ei} ومحددها، واستنتج منهما شرطي القضية~\ref{prop:ei}. على حدّ الشرط~(ii) يفقد الاتزان استقراره تذبذبيًا. أوجد دور هذه التذبذبات بأرقام الشكل~\ref{fig:ei}.
\end{exercise}

\begin{exercise}[\lvlB]
\label{ex:03:3}
\emph{إيقاع من التأخير.} لنستبدل بالتثبيط البطيء في النموذج~\eqref{eq:ei} تثبيطًا سريعًا بتأخير $d$: $\tau_E\dot E=-E(t)-GE(t-d)+h$. عند $\tau_E=10\ms$ أوجد أصغر تأخير $d_c$ يعطي تذبذبات، ودورها عند $G=2$ و$G=5$. هل تقع ضمن إيقاعات القشرة السريعة، $12\text{--}50\ms$، خلافًا للتمرين~\ref{ex:03:2}؟
\end{exercise}

\begin{exercise}[\lvlB]
\label{ex:03:4}
\emph{نمذجة: اختيار بذاكرة.} كامل~\eqref{eq:wta} عند $\beta=2$ و$\tau=1$. (أ)~عند $I_1=1$ ارفع $I_2$ ببطء من $0$ إلى $3$ ثم أعده: عند أي قيم $I_2$ تنقلب الشبكة؟ (ب)~بكم يزداد زمن القرار انطلاقًا من الصمت حين يصغر فرق المدخلين $\varepsilon$ عشر مرات؟
\end{exercise}

\begin{exercise}[\lvlB]
\label{ex:03:5}
\emph{تصميم: كاشف لنطاق من السرعات.} يجب أن يصمت كاشف الشكل~\ref{fig:direction} عند الحركة من $A$ إلى $B$ بزمن عبور $5\text{--}20\ms$. وسيط \nt{GABA} بتأخير $d_k$ يحظر $Y$ على المجال $[d_k,\,d_k+5\ms]$ بعد نبضة $A$؛ والإثارة من $B$ تصل فورًا. كم وسيطًا يلزم، وبأي تأخيرات؟
\end{exercise}

\begin{exercise}[\lvlB]
\label{ex:03:6}
\emph{الحظر والخلفية.} كم عصبونًا تتطلب الزوجية $x\oplus y\oplus z$ في المخطط العام ”وسيط \nt{GABA} لكل مدخل، وحظر لكل توليفة فيها $f=1$، و\foreignlanguage{english}{OR} مشتركة“؟ ابنها من عصبونين، \nt{GABA} و\nt{GLUT}، يتلقيان المداخل مباشرة، مع بيان الأوزان والعتبات.
\end{exercise}

\begin{exercise}[\lvlC]
\label{ex:03:7}
\emph{تصميم: الاختيار من مئة.} يجب أن تثبط كل واحدة من $100$ مجموعة \nt{GLUT} كل المجموعات الأخرى بالتساوي، لا نفسها. وسطاء \nt{GABA} خطيون تثيرهم المجموعات ويثبطونها؛ ويمكن للمجموعات أن يثير بعضها بعضًا، لا نفسها. ما أصغر عدد من الوسطاء؟ وماذا لو لم تكن ثمة وصلات إثارة مباشرة أصلًا؟
\end{exercise}

\begin{exercise}[\lvlC]
\label{ex:03:8}
\emph{بحث: متى يكون التثبيط متناقضًا.} في نموذج الشكل~\ref{fig:ei} ($w_{EI}=2$، $w_{IE}=1$) صارت دالة نقل مجموعة \nt{GLUT} هي $f(u)=[u]_+^2$، وبقيت مجموعة \nt{GABA} خطية. فوق أي مدخل $h_c$ تخفض الإضافة إلى مجموعة \nt{GABA} ترددها هي نفسها؟ تحقق بالنمذجة.
\end{exercise}
